% ============================================================================
% ARBEITSBLATT DS 8 - Gruppe B: Mi persona favorita (erweitert)
% ============================================================================

\documentclass[11pt,a4paper]{article}

\usepackage[utf8]{inputenc}
\usepackage[T1]{fontenc}
\usepackage[ngerman]{babel}
\usepackage[left=2cm,right=2cm,top=1.8cm,bottom=1.5cm]{geometry}
\usepackage{helvet}
\usepackage[table]{xcolor}
\usepackage{tabularx}
\usepackage{array}
\usepackage{enumitem}
\usepackage{tcolorbox}
\usepackage{fancyhdr}
\usepackage{lastpage}
\usepackage{amssymb}

\renewcommand{\familydefault}{\sfdefault}

\definecolor{spanisch}{HTML}{c41e3a}
\definecolor{grau}{HTML}{666666}
\definecolor{hellgrau}{HTML}{f5f5f5}
\definecolor{gruppeB}{HTML}{1565c0}

\pagestyle{fancy}
\fancyhf{}
\fancyhead[L]{\small\textcolor{grau}{Spanisch Spätbeginner -- Gruppe B}}
\fancyhead[R]{\small\textcolor{grau}{AB-08-B}}
\fancyfoot[C]{\small\thepage/\pageref{LastPage}}
\renewcommand{\headrulewidth}{0.5pt}

\newcommand{\luecke}[1]{\underline{\hspace{#1}}}
\newcommand{\punktebox}[1]{\hfill\fbox{\small \_\_/#1}}
\newcommand{\es}[1]{\textit{#1}}

\newtcolorbox{merkkasten}[1][]{
    colback=white,
    colframe=gruppeB,
    fonttitle=\bfseries\color{gruppeB},
    title=#1,
    boxrule=1.5pt,
    arc=0pt,
    left=5pt, right=5pt, top=5pt, bottom=5pt
}

\newtcolorbox{tippbox}{
    colback=hellgrau,
    colframe=grau,
    boxrule=0.5pt,
    arc=0pt,
    left=5pt, right=5pt, top=3pt, bottom=3pt
}

\newcounter{aufgabe}
\newcommand{\aufgabe}[2]{%
    \stepcounter{aufgabe}%
    \vspace{0.8em}%
    \noindent\textbf{\theaufgabe.} #1 \punktebox{#2}%
    \vspace{0.3em}%
}

\begin{document}

% ============================================================================
% KOPFBEREICH
% ============================================================================
\noindent
\begin{tabularx}{\textwidth}{@{}Xr@{}}
    {\Large\textbf{\textcolor{gruppeB}{Mi persona favorita}}} & Name: \luecke{5cm} \\[0.3em]
    {\small DS 8 -- Creatividad y Juego} & Datum: \luecke{3cm} \\
\end{tabularx}

\vspace{0.3em}
\hrule
\vspace{0.8em}

% ============================================================================
% MERKKASTEN: Adjektive zur Personenbeschreibung (erweitert)
% ============================================================================
\begin{merkkasten}[Kasten 1: Adjetivos para describir personas -- Nivel avanzado]
\vspace{0.2em}

\begin{tabularx}{\textwidth}{@{}|X|X|X|@{}}
\hline
\rowcolor{hellgrau}
\textbf{Positiv} & \textbf{Beruflich} & \textbf{Persönlichkeit} \\
\hline
simpático/-a, amable & trabajador/-a & introvertido/-a \\
inteligente, listo/-a & responsable & extrovertido/-a \\
divertido/-a, gracioso/-a & organizado/-a & optimista \\
paciente, tolerante & profesional & pesimista \\
cariñoso/-a (liebevoll) & \luecke{2.5cm} & ambicioso/-a \\
generoso/-a (großzügig) & \luecke{2.5cm} & \luecke{2.5cm} \\
\hline
\end{tabularx}

\vspace{0.3em}
\textbf{Steigerung:} muy + Adj., bastante + Adj., un poco + Adj., demasiado + Adj.
\vspace{0.2em}
\end{merkkasten}

\vspace{0.8em}

% ============================================================================
% MERKKASTEN 2: Erweiterte Strukturen
% ============================================================================
\begin{merkkasten}[Kasten 2: Strukturen für komplexe Personenbeschreibungen]
\vspace{0.2em}

\begin{tabularx}{\textwidth}{@{}|p{4cm}|X|@{}}
\hline
\rowcolor{hellgrau}
\textbf{Struktur} & \textbf{Beispiel} \\
\hline
\textbf{ser + Adjektiv} & Mi madre \textbf{es} una persona muy amable. \\[0.3em]
\hline
\textbf{parecer + Adj.} & \textbf{Parece} serio, pero es muy divertido. \\
(scheinen) & \\[0.3em]
\hline
\textbf{le gusta + Infinitiv} & \textbf{Le gusta} ayudar a los demás. \\[0.3em]
\hline
\textbf{siempre / nunca} & \textbf{Siempre} está de buen humor. \\[0.3em]
\hline
\end{tabularx}

\vspace{0.3em}
\end{merkkasten}

\vspace{1em}

% ============================================================================
% AUFGABEN
% ============================================================================

\aufgabe{Wortschatz-Staffel: Schreibe Adjektive nach Kategorien. (3 Min.)}{6}

\begin{tippbox}
\textbf{Ziel:} Mindestens 12 verschiedene Adjektive in den richtigen Kategorien.
\end{tippbox}

\vspace{0.3em}
\begin{tabularx}{\textwidth}{@{}|X|X|X|@{}}
\hline
\rowcolor{hellgrau}
\textbf{Positiv} & \textbf{Neutral/Negativ} & \textbf{Beruf/Charakter} \\
\hline
1. \luecke{2.2cm} & 1. \luecke{2.2cm} & 1. \luecke{2.2cm} \\[0.5em]
2. \luecke{2.2cm} & 2. \luecke{2.2cm} & 2. \luecke{2.2cm} \\[0.5em]
3. \luecke{2.2cm} & 3. \luecke{2.2cm} & 3. \luecke{2.2cm} \\[0.5em]
4. \luecke{2.2cm} & 4. \luecke{2.2cm} & 4. \luecke{2.2cm} \\[0.5em]
\hline
\end{tabularx}

\vspace{0.8em}

\aufgabe{Ergänze \es{ser}, \es{estar} oder \es{parecer}.}{6}

\begin{enumerate}[label=\alph*), leftmargin=2em]
    \item Mi jefe \luecke{2cm} muy estricto, pero en realidad \luecke{2cm} muy amable.
    \item Hoy mi madre \luecke{2cm} cansada porque ha trabajado mucho.
    \item Mis amigos \luecke{2cm} personas muy divertidas.
    \item Él \luecke{2cm} serio, pero cuando lo conoces, \luecke{2cm} simpático.
    \item ¿Cómo \luecke{2cm} tu profesor de español?
    \item Ella siempre \luecke{2cm} de buen humor.
\end{enumerate}

\vspace{0.8em}

\aufgabe{Schreibe komplexe Sätze mit Konnektoren und Steigerung.}{8}

\begin{tippbox}
\textbf{Verwende:} \es{aunque} (obwohl), \es{sin embargo} (jedoch), \es{por eso} (deshalb), \es{ya que} (da)
\end{tippbox}

\vspace{0.3em}
\begin{enumerate}[label=\alph*), leftmargin=2em]
    \item Meine Mutter scheint streng, aber sie ist sehr geduldig.

    \vspace{0.3cm}
    \luecke{13cm}

    \item Mein Freund ist sehr intelligent, deshalb hilft er mir immer.

    \vspace{0.3cm}
    \luecke{13cm}

    \item Obwohl mein Vater viel arbeitet, hat er immer Zeit für mich.

    \vspace{0.3cm}
    \luecke{13cm}

    \item Da meine Großmutter sehr weise ist, frage ich sie oft um Rat.

    \vspace{0.3cm}
    \luecke{13cm}
\end{enumerate}

\vspace{0.8em}

\aufgabe{Poster-Text: Schreibe einen ausführlichen Text über deine Lieblingsperson.}{10}

\begin{tippbox}
\textbf{Schreibe 8--10 Sätze.} Nutze verschiedene Strukturen, Konnektoren und Steigerungen!
\end{tippbox}

\vspace{0.5em}
\textbf{Inhalt:}
\begin{itemize}[leftmargin=1.5em, nosep]
    \item Wer ist die Person? (Name, Beziehung, Alter)
    \item Wie ist sie? (mind. 4 Adjektive mit Steigerung)
    \item Was macht sie gerne?
    \item Warum ist sie wichtig für dich?
    \item Eine Anekdote / ein Beispiel
\end{itemize}

\vspace{0.5em}
\begin{tabularx}{\textwidth}{@{}|X|@{}}
\hline
\\[0.3em]
\luecke{14cm} \\[0.7em]
\luecke{14cm} \\[0.7em]
\luecke{14cm} \\[0.7em]
\luecke{14cm} \\[0.7em]
\luecke{14cm} \\[0.7em]
\luecke{14cm} \\[0.7em]
\luecke{14cm} \\[0.7em]
\luecke{14cm} \\[0.5em]
\hline
\end{tabularx}

\vspace{0.8em}

\aufgabe{Escape-Room-Rätsel: Löse die Rätsel und finde den Code.}{5}

\textbf{Rätsel 1:} Finde 6 Adjektive in der Wortschlange. Die Anfangsbuchstaben ergeben ein Wort.

\vspace{0.3em}
{\small SIMPÁTICOINTELIGENTECREATIVOPACIENTERESPONSABLEDIVERTIDO}

\vspace{0.3em}
Adjektive: \luecke{8cm}

Anfangsbuchstaben: \luecke{2cm} \luecke{2cm} \luecke{2cm} \luecke{2cm} \luecke{2cm} \luecke{2cm}

Der 3. Buchstabe als Zahl (A=1, B=2, C=3...): \fbox{\phantom{XX}} = Code-Ziffer 1

\vspace{0.5em}
\textbf{Rätsel 2:} Zähle, wie oft \es{ser} vorkommt (nicht \es{estar}!):

\vspace{0.3em}
{\small 1. Mi hermano \textbf{es} simpático. 2. Hoy \textbf{estoy} cansado. 3. Ella \textbf{está} enferma. 4. Nosotros \textbf{somos} estudiantes. 5. ¿\textbf{Estás} contento?}

\vspace{0.3em}
Anzahl \es{ser}: \fbox{\phantom{XX}} = Code-Ziffer 2

\vspace{0.5em}
\textbf{Rätsel 3:} Ordne zu: A-médico, B-profesor, C-artista, D-bombero mit 1-creativo, 2-valiente, 3-paciente, 4-organizado

\vspace{0.3em}
A = \luecke{0.8cm}, B = \luecke{0.8cm}, C = \luecke{0.8cm}, D = \luecke{0.8cm}

Addiere die Zahlen für A und B: \luecke{1cm} + \luecke{1cm} = \fbox{\phantom{XX}} = Code-Ziffer 3

\vspace{0.5em}
\textbf{Gesamtcode:} \fbox{\phantom{XX}} -- \fbox{\phantom{XX}} -- \fbox{\phantom{XX}}

\vspace{1em}
\hrule
\vspace{0.3em}
\begin{flushright}
\textbf{Gesamt: \luecke{1cm} / 35 Punkte}
\end{flushright}

\end{document}
